\subsection{Prime fields}

The field of fractions of the ring Z of rational integers is called the field of rational numbers and is denoted by Q (I, p.~117). For every prime number p the quotient ring \(\mathbf{Z}/(p)\) is a finite \(^{1}\) field of p elements, denoted by \(F_{p}\) in the sequel. The field Q is infinite since it contains Z, and is therefore not isomorphic to any field \(F_{p}\) . If p and \(p'\) are distinct prime numbers, the fields \(F_{p}\) and \(F_{p'}\) have distinct cardinals and so are not isomorphic.

DEFINITION 1. --- A field is said to be prime if it is isomorphic either to Q or to one of the fields \(F_{p}\) .

Every subfield of Q contains the ring Z and hence the field of fractions Q of Z; every subring of \(F_{p}\) is necessarily equal to \(F_{p}\) . Therefore every subfield of a prime field is necessarily equal to it (cf.~Cor. 2 of Th. 1 below). Let P be a prime field and A a ring; if f and \(f'\) are two homomorphisms of P into A, then the set of \(x \in P\) such that \(f(x) = f'(x)\) is a subfield of P, hence by what has been said we must have \(f = f'\) . In particular, the only endomorphism of a prime field is the identity mapping.

THEOREM 1. --- Let A be a ring; suppose that there exists a subfield of A. Then A has a unique subfield P which is a prime field. Moreover, P is contained in the centre of A and in every subfield of A.

Let K be a subfield of A, C the centre of A, and put \(K' = K \cap C\) ; then \(K'\) is a subfield of A. Let f be the unique homomorphism of Z into A and p its kernel. Every subring of A, in particular \(K'\) , contains \(f(\mathbf{Z})\) ; hence the ideal p is prime (I, p.~116-117). If \(\mathfrak{p} = (0)\) , the homomorphism f of Z into \(K'\) is injective; hence it extends (I, p.~116) to an isomorphism \(\bar{f}\) of Q onto a subfield P of \(K'\) . If \(\mathfrak{p} \neq (0)\) , there is a strictly positive integer p such that \(\mathfrak{p} = (p)\) (I, p.~111); if we had p = ab with a \textgreater{} 1 and b \textgreater{} 1, this would mean \(a \notin p\) , \(b \notin p\) and \(ab \in p\) in contradiction with the fact that p is prime. The number p is therefore prime and by passage to quotients f defines an isomorphism of \(F_p = Z/p\) onto a subfield P of \(K'\) . In both cases P is a subfield of A contained in the centre C of A, and it is a prime field. Let L be a subfield of A; then \(P \cap L\) is a subfield of P, and since P is prime, we have \(P \cap L = P\) , whence \(P \subset L\) . If \(P'\) is a subfield of A and is a prime field, then by what has been said, \(P \subset P'\) , whence \(P = P'\) because \(P'\) is a prime field.

COROLLARY 1. --- Let K be a field. There exists a unique subfield of K which is a prime field, and this is the least subfield of K.

COROLLARY 2. --- For a field to be prime it is necessary and sufficient that it should contain no subfield other than itself.

